
\documentclass[12pt]{article}

%% Mathematics
%\usepackage{amsmath,amssymb,amsthm}

%% Physics
%\usepackage{physics}

%% Customization of captions in floating environments
%\usepackage{caption} % for \caption*

%% Citation improvement
%\usepackage{cite}

%% Hypertext marks
\usepackage[hidelinks]{hyperref} % `hyperref` configuration must go before `orcidlink` to avoid the "Option clash for package hyperref" error.

%% ORCID
\usepackage{orcidlink} % to insert a hyperlinked ORCiD logo

%% Page layout settings
%\textwidth 160mm
%\textheight 230mm
%\voffset = -20mm

\title{Title of the manuscript}

\author{R.V.~Romanik \orcidlink{0009-0008-2405-9038}
	\\ \small Yukhnovskii Institute for Condensed Matter Physics, NAS of Ukraine 
	\\ \small 1~Svientsitskii Street, 79011, Lviv, Ukraine 
	\\ \small romanik@icmp.lviv.ua}

%\date{} % to remove the date

\begin{document}
	
	\maketitle
	
	\begin{abstract}
		Abstract text. Pseudo-abstract. We investigate the thermodynamic behavior of a generalized many-particle system with competing interactions within a self-consistent theoretical framework. The model incorporates nonlinear coupling between microscopic degrees of freedom and macroscopic response functions, allowing us to examine the emergence of collective phenomena under varying external conditions.		
		\\[0.5cm]
		\textbf{Keywords:} Simple fluids, critical behavior.
	\end{abstract}
	
	\tableofcontents
	
	\section{Introduction}	
	Some introductory text with a citation
	\cite{RDKPS26a,RDKPS26b}, \cite{KD22}.
	
	\textbf{Pseudo-Introduction.} The collective behavior of interacting many-particle systems
	has attracted considerable attention in modern statistical
	physics due to its importance for understanding complex
	thermodynamic phenomena \cite{PLA99,HansenMcDonald13,Johnston14}.
	In particular, the interplay between microscopic interactions
	and macroscopic fluctuations may lead to nontrivial
	modifications of the equation of state and the emergence
	of anomalous phase behavior \cite{WM02,AW07}.
	Despite substantial theoretical progress, several aspects
	of these phenomena remain insufficiently understood.
	
	\section{Model}
	
	The partition function is
	\begin{equation}
		\Xi = \sum_{N=0}^{\infty}
		e^{\beta\mu N} Z_N.
	\end{equation}
	
	\section{Results}
	
	\section{Conclusions}
	
	% Bibliography
	\bibliographystyle{elsarticle-num}
	\bibliography{articles,books}
	
\end{document}
